

\subsection{Definitions}

When we say that the congregation is the body of Christ here on earth, we are no longer speaking in terms of polity. Instead, we are making a dogmatic claim about the nature of the church.

For our purposes we will use the following definitions:

\begin{itemize}
    \item \textbf{Dogmatics:} What we confess because Scripture teaches it.

    \item \textbf{Ecclesiology:} What we believe Scripture teaches about the nature of the congregation.

    \item \textbf{Polity:} How we organize and operate the congregation. 
\end{itemize}

These three terms are tightly coupled, as polity depends on ecclesiology and ecclesiology depends on dogmatics. 

As we read Sverdrup, Oftedal, and Evjen, we need to be very careful with these terms. This may sound silly, but I recommend you make three notecards, one for each term. As you read this material, select and hold the appropriate card in your hand. Do this for each point. Even better, read the passage again with one of the wrong cards and see how it changes your point of view.

Much of the Evjen controversy involves holding the correct card at the correct time. \textbf{If you are holding a dogmatics card and your opponent is holding a polity card, the argument is already lost.} As you read the remainder of this book, always have the cards in hand and ask which is being used in any given argument.    

Let's try a few examples:

\paragraph{Example 1: The Bible is the inerrant word of God} 

This is a dogmatic claim describing how we stand under Scripture. We will not understand everything, and there are sections we will not like. Yet we receive the Word as it is written and struggle together to find meaning and application in our lives.
 

\paragraph{Example 2: We are all sinners} 

This is a dogmatic claim that points the way to Jesus, from Genesis to the Cross. Without your sin and my sin, there would be no reason for Jesus to die. Our best works appear as filthy rags. Jesus nailed our sins to the Cross.

\paragraph{Example 3: All members of the congregation are called (have a vocation) to use their gifts to help each other} 

This ecclesiology is encoded into the latter half of the AFLC's principles, especially P9 and P12.


\paragraph{Example 4: The Congregation is the body of Christ}

This starts out as dogmatics but refuses to stay there. Jesus is not stationary. Instead, he gives spirit and life to the congregation. Order and mission (polity) are gifts to guide the movement of the congregational body.

\paragraph{Example 5: ``Silent Night'' and ``Amazing Grace'' sung a cappella}

Let's call this a local polity consideration with its heart grounded in ecclesiology. It recognizes that when the congregation sings together, the members must be able to hear one another. Active voices are more important than a wall of music, no matter how beautiful. As with the Lord's Prayer, the congregation must be able to hear itself as we confesses our sinful nature and the promise of the Cross.

\begin{quote}
Amazing grace (how sweet the sound)\\
that saved a wretch like me!\\
I once was lost, but now am found,\\
was blind, but now I see.\\
\end{quote}


\paragraph{Example 6: Lutherans are confessional}

This statement is hard to categorize. It is an operational statement (polity) reflected in the weekly church bulletin with the creeds rotating on a monthly basis. But it refuses to let go of the nature of the church (ecclesiology) because that is what Lutherans believe (dogmatics). Lutheran laypeople occasionally refer to the Augsburg Confession, but let's be honest, that's a hard read.    

 
\paragraph{Example 7: Principle 1: ``According to the Word of God, the congregation is the right form of the Kingdom of God on earth.''}

This synthetic statement is so complex that it requires it own section.


\subsection{Crossing Boundaries with P1}

The first principle is the keel for the entire AFLC. It refuses to categorically fit within our clean definitions. Instead, it has elements of all three:


\begin{itemize}
    \item \textbf{According to the Word of God:} This is a dogmatic biblical teaching.

    \item \textbf{the congregation is the right form:} This is a dogmatic assertion leading to an ecclesiological claim about the nature of the congregation.

    \item \textbf{of the Kingdom of God on earth:} This is an ecclesiological description of the visible congregation. This profound heavenly confession will not be constrained. By implication, is pulls at the remaining eleven principles. 
\end{itemize} 


Finally, we must recognize that P1 does not stand by itself. It carries the weight of the 11 other principles which cast it as a unifying polity statement. Taken together, P1 is a dogmatic authority statement describing the ecclesiological nature of the congregation, with a polity call to duty. 

The bottom line is that P1 is not reducible. We are not going to solve this with a pure biblical claim, a clear ecclesiological, or a clear polity-based statement. Instead, using today's terminology, the principles are a holistic statement describing the foundations upon which the AFLC is built.

This is why those notecards are so important, as we need to be very careful with our AFLC apologetics. The principles are like a stair-stepped claim with lived consequences. It goes something like this:

\begin{itemize}
   \item Authority is given in the Word.
   \item Authority shapes the visible congregation.
   \item The congregation has a mission and acts with spirit and life.
\end{itemize} 


% Factcheck: To Quote Evjen: ``dogmatics gives the gift; ethics/polity receives the gift as task.''



